% Compile with XeLaTeX (Overleaf: Menu -> Compiler -> XeLaTeX).
% XeLaTeX reads Unicode directly, so Greek stays Greek and English stays English.
\documentclass[11pt,a4paper]{article}

\usepackage{fontspec}
\usepackage{polyglossia}
\setdefaultlanguage{greek}
\setotherlanguage{english}

% Computer Modern Unicode: the classic LaTeX look, with both Greek and Latin glyphs
\setmainfont{cmunrm}[Extension=.otf, BoldFont=cmunbx, ItalicFont=cmunti, BoldItalicFont=cmunbi]
\setsansfont{cmunss}[Extension=.otf, BoldFont=cmunsx, ItalicFont=cmunsi, BoldItalicFont=cmunso]
\setmonofont{cmuntt}[Extension=.otf, BoldFont=cmuntb, ItalicFont=cmunit, BoldItalicFont=cmuntx]

\usepackage[margin=2.5cm]{geometry}
\usepackage{amsmath}
\usepackage{booktabs}
\usepackage{tabularx}
\usepackage{enumitem}
\usepackage{xcolor}
\usepackage{fancyvrb}
\usepackage[hidelinks,hypertexnames=false]{hyperref}

\setlist{itemsep=2pt}
\newcommand{\code}[1]{\texttt{#1}}
\fvset{frame=single, framesep=6pt, fontsize=\small}

\begin{document}

% ---------------- Title page ----------------
\begin{titlepage}
  \centering
  \vspace*{3cm}
  {\LARGE\bfseries Εφαρμογή Διαχείρισης Εκδηλώσεων και\\[4pt] Ηλεκτρονικών Κρατήσεων\par}
  \vspace{1.5cm}
  {\large Μάθημα: Ανάπτυξη Εφαρμογών Παγκόσμιου Ιστού\par}
  \vspace{2cm}
  {\large\bfseries Ομάδα\par}
  \vspace{0.4cm}
  {\large Μπόνης Ματθαίος -- 1115202300136\par}
  {\large Φραντζιάς Αντρέας -- 1115202300234\par}
  \vfill
  {\large Σεπτέμβριος 2026\par}
\end{titlepage}

\tableofcontents
\newpage

% ---------------- 1 ----------------
\section{Εισαγωγή}

Στην εργασία υλοποιήθηκε μία πλήρης web εφαρμογή για τη διαχείριση εκδηλώσεων και την
ηλεκτρονική κράτηση εισιτηρίων ή θέσεων. Η εφαρμογή ακολουθεί την αρχιτεκτονική
browser/server, όπου το frontend εκτελείται στον φυλλομετρητή και επικοινωνεί με το
backend μέσω REST API.

Ο βασικός στόχος ήταν να καλυφθούν οι ρόλοι και οι λειτουργίες που ζητούνται στην
εκφώνηση: επισκέπτης, συμμετέχων, διοργανωτής και διαχειριστής. Παράλληλα, υλοποιήθηκαν
μηχανισμοί αυθεντικοποίησης, εξουσιοδότησης, διαχείρισης χωρητικότητας, ανταλλαγής
μηνυμάτων, εξαγωγής δεδομένων και συστάσεων εκδηλώσεων.

Η υλοποίηση χωρίζεται σε δύο βασικά μέρη:
\begin{itemize}
  \item \textbf{Frontend:} εφαρμογή Angular για τη διεπαφή χρήστη.
  \item \textbf{Backend:} NestJS REST API με βάση δεδομένων PostgreSQL μέσω Prisma ORM.
\end{itemize}

% ---------------- 2 ----------------
\section{Αρχιτεκτονική Συστήματος}

Η εφαρμογή σχεδιάστηκε ως τυπική εφαρμογή παγκόσμιου ιστού με διαχωρισμό frontend και
backend. Τα δύο μέρη επικοινωνούν αποκλειστικά μέσω HTTPS αιτημάτων προς το REST API.

\subsection{Frontend}

Το frontend υλοποιήθηκε με Angular και TypeScript. Περιλαμβάνει σελίδες για
καλωσόρισμα, σύνδεση, εγγραφή, αναζήτηση εκδηλώσεων, προβολή λεπτομερειών, κράτηση,
διαχείριση εκδηλώσεων, μηνύματα και λειτουργίες διαχειριστή. Για τους χάρτες
χρησιμοποιήθηκε η βιβλιοθήκη Leaflet με OpenStreetMap, τόσο στις λεπτομέρειες
εκδήλωσης όσο και στην επιλογή τοποθεσίας κατά την εγγραφή και τη δημιουργία εκδήλωσης.

Η επικοινωνία με το backend γίνεται μέσα από ξεχωριστά services ανά περιοχή του API
(π.χ.\ \code{EventsApiService}, \code{BookingsApiService}, \code{MessagingApiService}),
ώστε κάθε URL να ορίζεται σε ένα μόνο σημείο. Ένας \code{AuthInterceptor} προσθέτει
αυτόματα το JWT token στην κεφαλίδα \code{Authorization} κάθε αιτήματος και, αν το
backend απαντήσει με \code{401}, αποσυνδέει τον χρήστη και τον μεταφέρει στη σελίδα
σύνδεσης. Οι σελίδες προστατεύονται επιπλέον με τα route guards \code{AuthGuard} και
\code{RoleGuard}.

\subsection{Backend}

Το backend υλοποιήθηκε με NestJS και παρέχει REST API. Είναι υπεύθυνο για την
επιχειρησιακή λογική, την πρόσβαση στη βάση δεδομένων, τους ελέγχους ρόλων, την έκδοση
JWT tokens, την εξαγωγή XML/JSON και τον αλγόριθμο συστάσεων. Τα δεδομένα εισόδου
ελέγχονται με DTOs και \code{class-validator}, ώστε λανθασμένα αιτήματα να
απορρίπτονται με \code{400}.

\subsection{Βάση Δεδομένων}

Η βάση δεδομένων είναι PostgreSQL. Η πρόσβαση γίνεται μέσω Prisma ORM, το οποίο
αντιστοιχίζει το σχεσιακό μοντέλο σε TypeScript αντικείμενα και διευκολύνει τις
συναλλαγές (transactions), τους περιορισμούς και τις συσχετίσεις.

\subsection{Ασφάλεια και HTTPS}

Όλες οι αλληλεπιδράσεις εκτελούνται μέσω HTTPS. Το script εκκίνησης δημιουργεί
self-signed certificates για τοπική χρήση, ώστε να καλυφθεί η απαίτηση για SSL/TLS κατά
την ανάπτυξη και την επίδειξη της εφαρμογής. Οι κωδικοί αποθηκεύονται ως hash με
bcrypt και δεν επιστρέφονται ποτέ από το API.

% ---------------- 3 ----------------
\section{Σχεδίαση Βάσης Δεδομένων}

Το σχεσιακό μοντέλο σχεδιάστηκε με βάση το DTD της εκφώνησης και τις επιπλέον
λειτουργίες της εφαρμογής. Τα βασικά μοντέλα είναι τα ακόλουθα:

\begin{center}
\begin{tabularx}{\textwidth}{@{}l X@{}}
  \toprule
  \textbf{Οντότητα} & \textbf{Περιγραφή} \\
  \midrule
  \code{User} & Αποθηκεύει στοιχεία χρήστη, ρόλους, κατάσταση έγκρισης, στοιχεία
    επικοινωνίας, διεύθυνση, ΑΦΜ και γεωγραφική θέση. \\
  \code{Event} & Περιγράφει την εκδήλωση με τίτλο, τύπο, χώρο, διεύθυνση, πόλη, χώρα,
    συντεταγμένες, ημερομηνίες, χωρητικότητα, κατάσταση και περιγραφή. \\
  \code{Category} & Επιτρέπει σε κάθε εκδήλωση να ανήκει σε μία ή περισσότερες
    κατηγορίες. \\
  \code{TicketType} & Περιγράφει τύπους εισιτηρίων με όνομα, τιμή, ποσότητα και
    διαθέσιμο υπόλοιπο. \\
  \code{Booking} & Αποθηκεύει τις κρατήσεις, τον συμμετέχοντα, το είδος εισιτηρίου, τον
    αριθμό εισιτηρίων, το κόστος και την κατάσταση κράτησης. \\
  \code{Message} & Υποστηρίζει τα εισερχόμενα, απεσταλμένα, μη αναγνωσμένα μηνύματα και
    τη διαγραφή ανά χρήστη. \\
  \code{EventView} & Καταγράφει προβολές εκδηλώσεων από χρήστες και χρησιμοποιείται
    στον αλγόριθμο συστάσεων. \\
  \code{Photo} & Αποθηκεύει προαιρετικές φωτογραφίες ή URLs φωτογραφιών για κάθε
    εκδήλωση. \\
  \bottomrule
\end{tabularx}
\end{center}

Οι καταστάσεις ορίζονται ως enums:
\begin{itemize}
  \item \code{UserStatus}: \code{PENDING}, \code{APPROVED}, \code{REJECTED}
  \item \code{EventStatus}: \code{DRAFT}, \code{PUBLISHED}, \code{COMPLETED}, \code{CANCELLED}
  \item \code{BookingStatus}: \code{PENDING}, \code{CONFIRMED}, \code{CANCELLED}
\end{itemize}

Η χωρητικότητα των εκδηλώσεων ελέγχεται στο backend. Δεν επιτρέπεται το άθροισμα των
εισιτηρίων να υπερβαίνει τη συνολική χωρητικότητα και δεν επιτρέπεται κράτηση όταν δεν
υπάρχουν διαθέσιμα εισιτήρια για τον ζητούμενο τύπο.

% ---------------- 4 ----------------
\section{Ρόλοι και Αυθεντικοποίηση}

Η εφαρμογή υποστηρίζει τους ρόλους που ζητούνται στην εκφώνηση:
\begin{itemize}
  \item \textbf{Επισκέπτης:} μπορεί να βλέπει και να αναζητά δημοσιευμένες εκδηλώσεις,
    χωρίς να κάνει κρατήσεις.
  \item \textbf{Συμμετέχων:} μπορεί να κάνει κρατήσεις σε ενεργές εκδηλώσεις, να βλέπει
    τις κρατήσεις του και να χρησιμοποιεί το σύστημα μηνυμάτων.
  \item \textbf{Διοργανωτής:} μπορεί να δημιουργεί, να επεξεργάζεται, να δημοσιεύει, να
    ακυρώνει και να διαχειρίζεται τις δικές του εκδηλώσεις.
  \item \textbf{Διαχειριστής:} μπορεί να εγκρίνει ή να απορρίπτει χρήστες και να εξάγει
    τα δεδομένα των εκδηλώσεων.
\end{itemize}

Η σύνδεση γίνεται με username και password. Μετά την επιτυχημένη σύνδεση, το backend
εκδίδει JWT token. Το token αποστέλλεται από το frontend σε κάθε αίτημα προς
προστατευμένα endpoints, και τα guards του NestJS (\code{JwtAuthGuard},
\code{RolesGuard}) ελέγχουν αν ο χρήστης έχει τον κατάλληλο ρόλο. Τα δημόσια endpoints
σημειώνονται ρητά με τον decorator \code{@Public()}.

Νέοι χρήστες που εγγράφονται δεν μπορούν να συνδεθούν αμέσως. Η κατάστασή τους είναι
\code{PENDING} μέχρι να εγκριθούν από τον διαχειριστή.

% ---------------- 5 ----------------
\section{Λειτουργίες Εφαρμογής}

\subsection{Καλωσόρισμα, Σύνδεση και Εγγραφή}

Η πρώτη σελίδα της εφαρμογής δίνει πρόσβαση σε σύνδεση και εγγραφή. Αφού συνδεθεί, ο
χρήστης βλέπει τους ρόλους του, συνδέσμους προς τις σελίδες που του αντιστοιχούν και
προτεινόμενες εκδηλώσεις.

Η εγγραφή γίνεται σε τρία βήματα (λογαριασμός, στοιχεία, τοποθεσία). Ζητούνται τα
στοιχεία της εκφώνησης: username, password με επιβεβαίωση, όνομα, επώνυμο, email,
τηλέφωνο, ΑΦΜ, διεύθυνση και γεωγραφική θέση. Η θέση επιλέγεται πάνω σε χάρτη, με κλικ
ή μετακίνηση του δείκτη. Κάθε βήμα ελέγχεται πριν τη μετάβαση στο επόμενο, ενώ ο
έλεγχος ταύτισης των δύο κωδικών γίνεται σε επίπεδο φόρμας.

\subsection{Διαχείριση Χρηστών}

Ο διαχειριστής δημιουργείται κατά το seeding της βάσης. Από την κονσόλα διαχείρισης
μπορεί να δει τις αιτήσεις σε αναμονή και να τις εγκρίνει ή να τις απορρίψει, να
αναζητήσει χρήστες και να τους φιλτράρει ανά ρόλο (διοργανωτές, συμμετέχοντες), καθώς
και να δει τα πλήρη στοιχεία κάθε χρήστη.

\subsection{Διαχείριση Εκδηλώσεων}

Οι διοργανωτές μπορούν να δημιουργούν νέες εκδηλώσεις με κατηγορίες, τύπους
εισιτηρίων, χωρητικότητα, τοποθεσία και περιγραφή. Μια νέα εκδήλωση ξεκινά ως
\code{DRAFT} και γίνεται ορατή μόνο όταν δημοσιευτεί. Οι βασικοί κανόνες ελέγχονται στο
backend:
\begin{itemize}
  \item μόνο εκδηλώσεις \code{DRAFT} με τουλάχιστον έναν τύπο εισιτηρίου μπορούν να
    δημοσιευτούν,
  \item μόνο δημοσιευμένες εκδηλώσεις μπορούν να ακυρωθούν,
  \item οι τύποι εισιτηρίων δεν αλλάζουν αφού υπάρξουν κρατήσεις,
  \item εκδήλωση με κρατήσεις δεν μπορεί να διαγραφεί.
\end{itemize}

Στην ίδια σελίδα ο διοργανωτής βλέπει τις κρατήσεις κάθε εκδήλωσής του.

\subsection{Αναζήτηση και Προβολή Εκδηλώσεων}

Οι χρήστες μπορούν να αναζητήσουν εκδηλώσεις με βάση κατηγορία, κείμενο, ημερομηνίες,
τιμή και πόλη, με σελιδοποίηση των αποτελεσμάτων. Στο αίτημα αποστέλλονται μόνο τα
φίλτρα που έχουν συμπληρωθεί. Η σελίδα λεπτομερειών παρουσιάζει τα πλήρη στοιχεία της
εκδήλωσης και χάρτη OpenStreetMap με βάση τις αποθηκευμένες συντεταγμένες. Κάθε
προβολή καταγράφεται και χρησιμοποιείται από τον αλγόριθμο συστάσεων.

\subsection{Κρατήσεις}

Η κράτηση απαιτεί επιβεβαίωση από τον χρήστη. Μετά την οριστική υποβολή, η κράτηση
αποθηκεύεται και ενημερώνονται τα διαθέσιμα εισιτήρια μέσα σε ένα transaction. Ο έλεγχος
γίνεται στο backend, ώστε να μην μπορεί να παρακαμφθεί από το frontend. Οι κρατήσεις
είναι οριστικές και ο συμμετέχων τις βλέπει στη σελίδα «Οι κρατήσεις μου».

\subsection{Μηνύματα}

Υλοποιήθηκε σελίδα μηνυμάτων με εισερχόμενα, απεσταλμένα, κατάσταση ανάγνωσης,
απάντηση και διαγραφή. Η γραμμή πλοήγησης εμφανίζει τον αριθμό των μη αναγνωσμένων
μηνυμάτων, ο οποίος ανανεώνεται κάθε 30 δευτερόλεπτα (polling), αφού το HTTP δεν
επιτρέπει στον server να στείλει ειδοποίηση από μόνος του.

Κατά την ακύρωση εκδήλωσης, ο διοργανωτής γράφει ένα μήνυμα που αποστέλλεται σε όλους
τους συμμετέχοντες με κράτηση, και στη συνέχεια η εκδήλωση ακυρώνεται («Cancel +
notify»).

\subsection{Εξαγωγή XML και JSON}

Ο διαχειριστής μπορεί να εξάγει όλες τις εκδηλώσεις σε XML και JSON από την καρτέλα
«Data Export» της κονσόλας διαχείρισης. Η έξοδος XML ακολουθεί τη δομή του DTD της
εκφώνησης, με στοιχεία όπως \code{Event}, \code{Category}, \code{TicketTypes},
\code{Bookings}, \code{Organizer} και \code{Media}.

Τα endpoints εξαγωγής (\code{/export/events.xml} και \code{/export/events.json})
επιτρέπονται μόνο σε διαχειριστές. Γι' αυτό η λήψη γίνεται μέσω του \code{HttpClient} και όχι με απλό
σύνδεσμο: έτσι ο \code{AuthInterceptor} προσθέτει το token και το αρχείο αποθηκεύεται
τοπικά ως \code{events.xml} ή \code{events.json}.

% ---------------- 6 ----------------
\section{Αλγόριθμος Συστάσεων}

Ο αλγόριθμος συστάσεων υλοποιήθηκε από την αρχή, χωρίς έτοιμη βιβλιοθήκη
\textenglish{machine learning}. Χρησιμοποιήθηκε \textbf{\textenglish{Biased Matrix Factorization}}, όπου κάθε χρήστης και
κάθε εκδήλωση αναπαρίστανται με λανθάνοντα διανύσματα. Η πρόβλεψη ενδιαφέροντος του
χρήστη $u$ για την εκδήλωση $i$ δίνεται από τη σχέση:
\[
  \hat{r}_{ui} = \mu + b_u + b_i + \mathbf{p}_u^{\top}\mathbf{q}_i
\]
όπου $\mu$ είναι ο συνολικός μέσος όρος, $b_u$ και $b_i$ οι αποκλίσεις (bias) χρήστη και
εκδήλωσης, και $\mathbf{p}_u$, $\mathbf{q}_i$ τα λανθάνοντα διανύσματα.

Η εκπαίδευση γίνεται με stochastic gradient descent (learning rate $0.01$,
regularization $0.02$, $50$ epochs). Ως σήματα ενδιαφέροντος χρησιμοποιούνται:
\begin{itemize}
  \item επιβεβαιωμένες κρατήσεις, με βάρος $1.0$,
  \item προβολές εκδηλώσεων, με βάρος $0.3$ ανά προβολή και μέγιστο $0.8$,
  \item αρνητικά δείγματα (βάρος $0$) για εκδηλώσεις με τις οποίες ο χρήστης δεν
    αλληλεπίδρασε.
\end{itemize}

Αν ένας χρήστης δεν έχει ιστορικό, η εφαρμογή επιστρέφει προτάσεις με βάση τη
δημοφιλία των εκδηλώσεων (popularity fallback). Με αυτόν τον τρόπο εμφανίζονται
προτάσεις ακόμη και όταν τα δεδομένα είναι περιορισμένα. Λόγω του μικρού μεγέθους των
δεδομένων, το μοντέλο εκπαιδεύεται εκ νέου σε κάθε αίτημα ώστε το αποτέλεσμα να είναι
πάντα ενημερωμένο. Για μεγάλα δεδομένα θα χρειαζόταν περιοδική εκπαίδευση.

Ο βασικός κώδικας βρίσκεται στα αρχεία:
\begin{itemize}
  \item \code{backend/src/recommendations/bmf.ts}
  \item \code{backend/src/recommendations/recommendations.service.ts}
\end{itemize}

% ---------------- 7 ----------------
\section{Εγκατάσταση και Εκτέλεση}

Για την εκτέλεση χρειάζονται Node.js, npm, Docker και ένας σύγχρονος browser. Η πιο
απλή εκκίνηση γίνεται από τον ριζικό φάκελο του project:
\begin{Verbatim}
./run-all.sh
\end{Verbatim}

Το script εκτελεί τα βασικά βήματα αυτόματα και μπορεί να εκτελεστεί ξανά με ασφάλεια:
\begin{itemize}
  \item δημιουργεί ή ξεκινά PostgreSQL container,
  \item δημιουργεί τοπικό \code{backend/.env} αν λείπει,
  \item δημιουργεί self-signed TLS certificates,
  \item εγκαθιστά τα dependencies,
  \item εκτελεί τα Prisma migrations,
  \item δημιουργεί demo χρήστες και demo εκδήλωση (seed),
  \item ξεκινά backend και frontend.
\end{itemize}

Με \code{Ctrl+C} σταματούν backend και frontend, ενώ το container της βάσης παραμένει
ενεργό ώστε να μη χάνονται τα δεδομένα. Οι διευθύνσεις της εφαρμογής είναι:
\begin{itemize}
  \item Backend: \code{https://localhost:3000}
  \item Frontend: \code{https://localhost:4200}
\end{itemize}

Οι demo λογαριασμοί είναι:
\begin{center}
\begin{tabular}{@{}lll@{}}
  \toprule
  \textbf{Username} & \textbf{Password} & \textbf{Ρόλος} \\
  \midrule
  \code{admin}     & \code{admin123}     & Διαχειριστής \\
  \code{organizer} & \code{organizer123} & Διοργανωτής \\
  \code{user}      & \code{user1234}     & Συμμετέχων \\
  \bottomrule
\end{tabular}
\end{center}

% ---------------- 8 ----------------
\section{Έλεγχος και Δοκιμές}

Για τον έλεγχο του backend χρησιμοποιήθηκαν unit tests. Κατά την τελική επαλήθευση
πέρασαν και τα 14 test suites (36 tests):
\begin{Verbatim}
cd backend
npm test
\end{Verbatim}

Επίσης έγινε build του frontend:
\begin{Verbatim}
cd frontend
npm run build
\end{Verbatim}

Το build ολοκληρώθηκε επιτυχώς. Εμφανίστηκαν μόνο δύο προειδοποιήσεις της Angular: ότι
το αρχικό bundle ξεπερνά ελάχιστα το όριο μεγέθους (περίπου 906~kB έναντι 900~kB) και
ότι η Leaflet είναι CommonJS dependency. Οι προειδοποιήσεις αυτές δεν εμποδίζουν τη
λειτουργία της εφαρμογής.

Πέρα από τα αυτοματοποιημένα checks, έγινε χειροκίνητος έλεγχος των βασικών ροών:
εγγραφή, έγκριση χρήστη, σύνδεση, δημιουργία εκδήλωσης, αναζήτηση, προβολή χάρτη,
κράτηση, ανταλλαγή μηνυμάτων, εξαγωγή αρχείων και εμφάνιση συστάσεων.

% ---------------- 9 ----------------
\section{Παραδοχές και Σχεδιαστικές Αποφάσεις}

Κατά την υλοποίηση έγιναν οι παρακάτω παραδοχές:
\begin{itemize}
  \item Τα TLS certificates είναι self-signed, επειδή η εφαρμογή εκτελείται τοπικά για
    σκοπούς ανάπτυξης και αξιολόγησης.
  \item Η βάση δεδομένων εκτελείται σε Docker container για να είναι εύκολη η
    εγκατάσταση.
  \item Το αρχείο \code{.env} δεν παραδίδεται, επειδή περιέχει τοπικές ρυθμίσεις. Το
    \code{run-all.sh} το δημιουργεί αυτόματα για τοπική εκτέλεση.
  \item Τα \code{node\_modules}, τα build artifacts και τα παραγόμενα certificates δεν
    περιλαμβάνονται στο zip του πηγαίου κώδικα.
  \item Το dataset εισαγωγής δεν περιλαμβάνεται στο zip, εφόσον παρέχεται ξεχωριστά από
    το μάθημα.
  \item Οι κρατήσεις είναι οριστικές και δεν ακυρώνονται από τον συμμετέχοντα. Όταν
    ακυρώνεται μια εκδήλωση, οι υπάρχουσες κρατήσεις διατηρούνται και οι συμμετέχοντες
    ενημερώνονται με μήνυμα.
\end{itemize}

% ---------------- 10 ----------------
\section{Δυσκολίες}

Μία βασική δυσκολία ήταν η μετατροπή της περιγραφής XML/DTD σε σωστό σχεσιακό μοντέλο.
Η λύση ήταν να χωριστούν οι επαναλαμβανόμενες πληροφορίες σε ξεχωριστούς πίνακες, όπως
\code{TicketType}, \code{Booking}, \code{Category} και \code{Photo}.

Άλλη σημαντική δυσκολία ήταν η διατήρηση της σωστής διαθεσιμότητας εισιτηρίων. Για τον
λόγο αυτό οι έλεγχοι γίνονται στο backend μέσα σε transactions και όχι μόνο στο frontend.

Η εξαγωγή αρχείων απαιτούσε επίσης προσοχή. Επειδή τα endpoints επιτρέπονται μόνο σε
διαχειριστές, ένας απλός σύνδεσμος λήψης απέτυχε με \code{401}, αφού ο browser δεν
στέλνει το JWT token. Η λύση ήταν η λήψη μέσω \code{HttpClient}, ώστε το token να
προστίθεται από τον interceptor.

Τέλος, ο αλγόριθμος συστάσεων απαιτούσε προσεκτικό χειρισμό των περιορισμένων
δεδομένων. Γι' αυτό χρησιμοποιήθηκαν τόσο κρατήσεις όσο και προβολές εκδηλώσεων, καθώς
και popularity fallback όταν δεν υπάρχει ιστορικό χρήστη.

% ---------------- 11 ----------------
\section{Συμπέρασμα}

Η εργασία ολοκληρώθηκε ως πλήρης εφαρμογή διαχείρισης εκδηλώσεων και ηλεκτρονικών
κρατήσεων. Υλοποιήθηκαν οι βασικές απαιτήσεις της εκφώνησης, όπως ρόλοι χρηστών,
έγκριση εγγραφών, CRUD εκδηλώσεων, έλεγχος χωρητικότητας, κρατήσεις, messaging, εξαγωγή
XML/JSON και σύστημα συστάσεων.

Μέσα από την υλοποίηση έγινε πρακτική εφαρμογή τεχνολογιών web, σχεσιακής βάσης
δεδομένων, REST API, JWT authentication και ενός απλού recommender system. Μελλοντικές
βελτιώσεις θα μπορούσαν να περιλαμβάνουν αναλυτικότερο UI για στατιστικά διοργανωτών
(π.χ.\ έσοδα ανά εκδήλωση), καλύτερη αναζήτηση με full-text indexes, ειδοποιήσεις σε
πραγματικό χρόνο μέσω WebSockets και εκτενέστερη αξιολόγηση του αλγορίθμου συστάσεων με
μεγαλύτερο dataset.

\end{document}
